\subsection{BÀI TẬP TRẮC NGHIỆM TỰ LUYỆN}
\subsection*{\indam{PHẦN I. Câu trắc nghiệm nhiều phương án lựa chọn. Thí sinh trả lời từ câu 1 đến câu 12. Mỗi câu hỏi thí sinh chỉ chọn một phương án.}} 
	\setcounter{ex}{0}
	\Opensolutionfile{ans}[ans/B4-De2-1]

\begin{ex}%[2H5H3-1]
	Trong không gian $Oxyz$, điểm nào sau đây nằm trong mặt cầu $(S)\colon \left(x-1\right)^2+\left(y-4\right)^2+\left(z-3\right)^2=16$?
	\choice
	{$M\left(0;7;-3\right)$}
	{$P\left(1;0;0\right)$}
	{\True $N\left(0;4;3\right)$}
	{$Q\left(1;0;3\right)$}
	\loigiai{
		Mặt cầu $(S)$ có tâm $I\left(1;4;3\right)$ và bán kính $R=4$.\\
		$IM=\sqrt{(-1)^2+3^2+(-6)^2}=\sqrt{46} > R\Rightarrow M$ nằm ngoài mặt cầu.\\
		$IN=\sqrt{(-1)^2+0^2+0^2}=1< R\Rightarrow N$ nằm trong mặt cầu.\\
		$IP=\sqrt{0^2+\left(-4\right)^2+\left(-3\right)^2}=5> R\Rightarrow P$ nằm ngoài mặt cầu.\\
		$IQ=\sqrt{0^2+\left(-4\right)^2+0^2}=4=R\Rightarrow Q$ thuộc mặt cầu.\\
	}
\end{ex}

%G:\My Drive\CODE12-2024\DE-ON-THEO BAI\2H5-TACH DE\Bai4-De2.tex
\begin{ex}%[2H5H3-3]
	Trong không gian $Oxyz$, cho mặt cầu $(S)$ có tâm $A\left(3;-1;1\right)$ và đi qua $M\left(2;-2;4\right)$. Phương trình mặt cầu $(S)$ là
	\choice
	{\True $\left(x-3\right)^2+\left(y+1\right)^2+\left(z-1\right)^2=11$}
	{$\left(x+3\right)^2+\left(y-1\right)^2+\left(z+1\right)^2=11$}
	{$\left(x+3\right)^2+\left(y-1\right)^2+\left(z+1\right)^2=\sqrt{11}$}
	{$\left(x-3\right)^2+\left(y+1\right)^2+\left(z-1\right)^2=\sqrt{11}$}
	\loigiai{
		Bán kính $R=AM=\sqrt{(-1)^2+(-1)^2+3^2}=\sqrt{11}$.\\
		Phương trình mặt cầu  là  $(S)\colon \left(x-3\right)^2+\left(y+1\right)^2+\left(z-1\right)^2=11$.
	}
\end{ex}

%G:\My Drive\CODE12-2024\DE-ON-THEO BAI\2H5-TACH DE\Bai4-De2.tex
\begin{ex}%[2H5N3-2]
	Trong không gian $Oxyz$, cho mặt cầu $(S):\left(x-2\right)^2+\left(y+2\right)^2+\left(z-1\right)^2=18$. Bán kính của $(S)$ bằng
	\choice
	{$9$}
	{$18$}
	{$6\sqrt{2}$}
	{\True $3\sqrt{2}$}
	\loigiai{
		Bán kính $R=\sqrt{18}=3\sqrt{2}$.
	}
\end{ex}

%G:\My Drive\CODE12-2024\DE-ON-THEO BAI\2H5-TACH DE\Bai4-De2.tex
\begin{ex}%[2H5H3-3]
	Trong không gian $Oxyz$, cho hai điểm $M\left(3;-2;5\right)$, $N\left(-1;6;-3\right)$. Mặt cầu đường kính $MN$ có phương trình là
	\choice
	{$\left(x+1\right)^2+\left(y+2\right)^2+\left(z+1\right)^2=6$}
	{$\left(x-1\right)^2+\left(y-2\right)^2+\left(z-1\right)^2=6$}
	{\True $\left(x-1\right)^2+\left(y-2\right)^2+\left(z-1\right)^2=36$}
	{$\left(x+1\right)^2+\left(y+2\right)^2+\left(z+1\right)^2=36$}
	\loigiai{
		Tâm $I$ của mặt cầu là trung điểm đoạn $MN$ $\Rightarrow$ $I\left(1;2;1\right)$.\\
		Bán kính mặt cầu $R=\dfrac{MN}{2}=\dfrac{\sqrt{\left(-1-3\right)^2+\left(6+2\right)^2+\left(-3-5\right)^2}}{2}=6$.\\
		Vậy phương trình mặt cầu là $\left(x-1\right)^2+\left(y-2\right)^2+\left(z-1\right)^2=36$.
	}
\end{ex}

%G:\My Drive\CODE12-2024\DE-ON-THEO BAI\2H5-TACH DE\Bai4-De2.tex
\begin{ex}%[2H5H3-2]
	Trong không gian $Oxyz$, cho mặt cầu $(S)\colon \left(x-5\right)^2+\left(y-1\right)^2+\left(z+2\right)^2=9$. Đường kính của mặt cầu $(S)$ là
	\choice
	{$9$}
	{$3$}
	{\True $6$}
	{$18$}
	\loigiai{
		Bán kính $R=\sqrt{9}=3\Rightarrow$ đường kính bằng $2R=2\cdot 3=6$.
	}
\end{ex}

%G:\My Drive\CODE12-2024\DE-ON-THEO BAI\2H5-TACH DE\Bai4-De2.tex
\begin{ex}%[2H5H3-2]
	Trong không gian $Oxyz$, phương trình nào sau đây là phương trình của mặt cầu?
	\choice
	{$x^2+z^2+3x-2y+4z-1=0$}
	{$x^2+y^2+z^2-2x+2y-4z+8=0$}
	{\True $x^2+y^2+z^2-2x+4z-1=0$}
	{$x^2+y^2+z^2+2xy-4y+4z-1=0$}
	\loigiai{
		Điều kiện để phương trình $x^2+y^2+z^2-2ax-2by-2cz+d=0$ là phương trình mặt cầu là
		\[
		a^2+b^2+c^2-d > 0.
		\]
		Xét phương trình  $x^2+y^2+z^2-2x+4z-1=0$ có  $1^2+0^2+(-2)^2-(-1)=6> 0$. \\
		Suy ra phương trình  $x^2+y^2+z^2-2x+4z-1=0$  là phương trình của một mặt cầu.
	}
\end{ex}

%G:\My Drive\CODE12-2024\DE-ON-THEO BAI\2H5-TACH DE\Bai4-De2.tex
\begin{ex}%[2H5H3-2]
	Trong không gian với hệ toạ độ $Oxyz$,  cho mặt cầu $(S):(x+2)^2+y^2+\left(z-3\right)^2=4$. Tâm của $(S)$ có toạ độ là
	\choice
	{$\left(-1;0;\dfrac{3}{2} \right)$}
	{$\left(2;0;-3\right)$}
	{$\left(1;0;\dfrac{3}{2} \right)$}
	{\True $\left(-2;0;3\right)$}
	\loigiai{
		Phương trình mặt cầu $\left(x-a\right)^2+\left(y-b\right)^2+\left(z-c\right)^2=R^2$ có tâm $I\ (a;b;c)$.\\
		Suy ra $(S)\colon (x+2)^2+y^2+\left(z-3\right)^2=4$ có tâm $\left(-2;0;3\right)$.
	}
\end{ex}

%G:\My Drive\CODE12-2024\DE-ON-THEO BAI\2H5-TACH DE\Bai4-De2.tex
\begin{ex}%[2H5H3-4]
	Trong không gian $Oxyz$, một thiết bị phát sóng đặt tại vị trí $A(2;0;0)$. Vùng phủ sóng của thiết bị có bán kính bằng $1$. Điểm nào sau đây thuộc vùng phủ sóng của thiết bị nói trên?
	\choice
	{\True $P\left(1;0;0\right)$}
	{$O\left(0;0;0\right)$}
	{$N\left(0;1;1\right)$}
	{$M\left(1;0;3\right)$}
	\loigiai{
		Ta có $AM=\sqrt{(-1)^2+0^2+3^2}=\sqrt{10} > R$ nên $M$ không thuộc vùng phủ sóng.\\
		$AN=\sqrt{0^2+1^2+1^2}=\sqrt{2} > R$ nên $N$ không thuộc vùng phủ sóng.\\
		$AP=\sqrt{(-1)^2+0^2+0^2}=1=R$ nên $P$ thuộc vùng phủ sóng.\\
		$AO=\sqrt{(-2)^2+0^2+0^2}=2> R$ nên $O$ không thuộc vùng phủ sóng.
	}
\end{ex}

%G:\My Drive\CODE12-2024\DE-ON-THEO BAI\2H5-TACH DE\Bai4-De2.tex
\begin{ex}%[2H5H3-2]
	Trong không gian $Oxyz$, cho mặt cầu $(S)\colon x^2+y^2+z^2-2x+4y+2z-3=0$. Bán kính $R$ của mặt cầu $(S)$ bằng
	\choice
	{$\sqrt{3}$}
	{$9$}
	{$3\sqrt{3}$}
	{\True $3$}
	\loigiai{
		Mặt cầu có dạng $x^2+y^2+z^2-2ax-2by-2cz+d=0$ có bán kính là $R=\sqrt{a^2+b^2+c^2-d}$.\\
		Suy ra $R=\sqrt{1^2+(-2)^2+1^2-(-3)}=3$.
	}
\end{ex}

%G:\My Drive\CODE12-2024\DE-ON-THEO BAI\2H5-TACH DE\Bai4-De2.tex
\begin{ex}%[2H5H3-3]
	Trong không gian $Oxyz$, phương trình mặt cầu tâm $I\left(1;2;3\right)$, bán kính $R=2$ có dạng
	\choice
	{\True $\left(x-1\right)^2+\left(y-2\right)^2+\left(z-3\right)^2=4$}
	{$\left(x-1\right)^2+\left(y-2\right)^2+\left(z-3\right)^2=2$}
	{$\left(x+1\right)^2+\left(y+2\right)^2+\left(z+3\right)^2=2$}
	{$\left(x+1\right)^2+\left(y+2\right)^2+\left(z+3\right)^2=4$}
	\loigiai{
		Phương trình mặt cầu tâm $I\left(1;2;3\right)$, bán kính $R=2$ có dạng
		\[
		\left(x-1\right)^2+\left(y-2\right)^2+\left(z-3\right)^2=4.
		\]
	}
\end{ex}

%G:\My Drive\CODE12-2024\DE-ON-THEO BAI\2H5-TACH DE\Bai4-De2.tex
\begin{ex}%[2H5H3-2]
	Trong không gian $Oxyz$, tìm $m$ để phương trình $x^2+y^2+z^2-2x-y+4z-m=0$ là phương trình của mặt cầu.
	\choice
	{$m\le \dfrac{21}{4}$}
	{\True $m > \dfrac{21}{4}$}
	{$m\ge \dfrac{21}{4}$}
	{$m < \dfrac{21}{4}$}
	\loigiai{
		Phương trình đã cho là phương trình mặt cầu khi và chỉ khi
		\[
		a^2+b^2+c^2-d > 0 \Leftrightarrow 1^2+\left(\dfrac{1}{2} \right)^2+(-2)^2+m > 0 \Leftrightarrow m > \dfrac{21}{4}.
		\]
	}
\end{ex}

%G:\My Drive\CODE12-2024\DE-ON-THEO BAI\2H5-TACH DE\Bai4-De2.tex
\begin{ex}%[2H5N3-1]
	Trong không gian với hệ toạ độ $Oxyz$, mặt cầu  $(S)\colon x^2+y^2+z^2-2ax-2by-2cz+d=0$ có bán kính $R$ bằng
	\choice
	{$a^2+b^2+c^2+d$}
	{$\sqrt{a^2+b^2+c^2+d}$}
	{\True $\sqrt{a^2+b^2+c^2-d}$}
	{$a^2+b^2+c^2-d$}
	\loigiai{
		Mặt cầu  $(S)\colon x^2+y^2+z^2-2ax-2by-2cz+d=0$ có bán kính $R = \sqrt{a^2+b^2+c^2-d}$.
	}
\end{ex}
	\Closesolutionfile{ans}

\subsection*{\indam{PHẦN II. Câu trắc nghiệm đúng sai. Thí sinh trả lời từ câu 1 đến câu 4. Trong mỗi ý a), b), c), d) ở mỗi câu, thí sinh chọn đúng hoặc sai.}}
	\setcounter{ex}{0}
	\Opensolutionfile{ans}[ans/B4-De2-2]
\begin{ex}%[2H5H3-3]
	Trong không gian $Oxyz$, cho mặt cầu $(S): \left(x-2\right)^2+y^2+\left(z+1\right)^2=1$ và mặt phẳng $(P)\colon  x+2y-z+1=0$. Các mệnh đề sau đúng hay sai?
	\choiceTF
	{Khoảng cách từ tâm $I$ đến mặt phẳng $(P)$ bằng $\dfrac{\sqrt{6}}{3}$}
	{\True Mặt cầu $(S)$ có tâm $I\left(2;0;-1\right)$ và bán kính $R=1$}
	{Mặt phẳng $(P)$ tiếp xúc mặt cầu $(S)$}
	{\True Phương trình mặt cầu tâm $I\left(2;0;-1\right)$ và tiếp xúc mặt phẳng $(P)$ là: $\left(S'\right): x^2+y^2+z^2-4x+2z+\dfrac{7}{3}=0$}
	\loigiai{
		\begin{itemchoice}
			\itemch \textbf{Sai}.\\
			$\mathrm{d}\left(I,(P)\right)=\dfrac{\left|2+1+1\right|}{\sqrt{1+4+1}}=\dfrac{2\sqrt{6}}{3}$.
			\itemch \textbf{Đúng}.\\
			Mặt cầu $(S)\colon \left(x-2\right)^2+y^2+\left(z+1\right)^2=1$ có tâm $I\left(2;0;-1\right)$ và bán kính $R=1$.
			\itemch \textbf{Sai}.\\
			Vì $\mathrm{d}\left(I,(P)\right)=\dfrac{2\sqrt{6}}{3} > R$ nên $(P)$ không cắt mặt cầu $(S)$.
			\itemch \textbf{Đúng}.\\
			Phương trình mặt cầu cần tìm có bán kính $R=\mathrm{d}\left(I,(P)\right)=\dfrac{2\sqrt{6}}{3}$, tâm $I\left(2;0;-1\right)$ nên có phương trình 
			\[
			\left(x-2\right)^2+y^2+\left(z+1\right)^2=\left(\dfrac{2\sqrt{6}}{3} \right)^2 \Leftrightarrow x^2+y^2+z^2-4x+2z+\dfrac{7}{3}=0.
			\]
		\end{itemchoice}
	}
\end{ex}

%G:\My Drive\CODE12-2024\DE-ON-THEO BAI\2H5-TACH DE\Bai4-De2.tex
\begin{ex}%[2H5H3-3]
	Trong không gian $Oxyz$, cho mặt cầu $(S)\colon \left(x-1\right)^2+(y+3)^2+\left(z-2\right)^2=49$. 
	\choiceTF
	{\True Mặt cầu $(S)$ có bán kính $R=7$}
	{\True Điểm $A\left(1;4;2\right)$ nằm trên mặt cầu $(S)$}
	{Mặt cầu $(S)$ có tâm $I\left(1;3;2\right)$}
	{Mặt cầu $(S)$ còn có phương trình: $x^2+y^2+z^2-2x+6y-4z-49=0$}
	\loigiai{
		\begin{itemchoice}
			\itemch \textbf{Đúng}.\\
			Mặt cầu $(S)\colon \left(x-1\right)^2+(y+3)^2+\left(z-2\right)^2=49=7^2$ nên mặt cầu $(S)$ có bán kính $R=7$.
			\itemch \textbf{Đúng}.\\
			Thay toạ độ điểm $A\left(1;4;2\right)$ vào vế trái của mặt cầu $(S)$ được: $VT=7^2=49$.
			\itemch \textbf{Sai}.\\
			Mặt cầu $(S)$ có tâm là $I\left(1;-3;2\right)$.
			\itemch \textbf{Sai}.\\
			Ta có
			\begin{eqnarray*}
				& & (S) \colon \left(x-1\right)^2+(y+3)^2+\left(z-2\right)^2=49\\
				& \Leftrightarrow &  x^2+y^2+z^2-2x+6y-4z+1+9+4=49\\
				& \Leftrightarrow & x^2+y^2+z^2-2x+6y-4z-35=0.
			\end{eqnarray*}
		\end{itemchoice}
	}
\end{ex}

%G:\My Drive\CODE12-2024\DE-ON-THEO BAI\2H5-TACH DE\Bai4-De2.tex
\begin{ex}%[2H5H3-3]
	Trong không gian $Oxyz$, cho mặt cầu $(S): x^2+y^2+z^2+2x+8y+1=0$. Các mệnh đề sau đúng hay sai?
	\choiceTF
	{Mặt cầu $(S)$ có tâm $I\left(1;4;0\right)$}
	{\True Mặt cầu $(S)$ còn có phương trình: $(S): \left(x+1\right)^2+\left(y+4\right)^2+z^2=16$}
	{\True Điểm $M\left(0;3;4\right)$ nằm bên ngoài mặt cầu $(S)$}
	{\True Mặt cầu $(S)$ có bán kính $R=4$}
	\loigiai{
		\begin{itemchoice}
			\itemch \textbf{Sai}.\\
			Mặt cầu $(S)$ có tâm là $I\left(-1;-4;0\right)$.
			\itemch \textbf{Đúng}.\\
			Vì mặt cầu $(S)$ có tâm $I\left(-1;-4;0\right)$, bán kính $R=4$ nên có phương trình là
			\[
			(S)\colon  \left(x+1\right)^2+\left(y+4\right)^2+z^2=16.
			\]
			\itemch \textbf{Đúng}.\\
			Ta có $IM=\sqrt{1^2+7^2+4^2} > R$ nên điểm $M\left(0;3;4\right)$ nằm bên ngoài mặt cầu $(S)$.
			\itemch \textbf{Đúng}.\\
			Mặt cầu $(S)$ có bán kính $R=\sqrt{a^2+b^2+c^2-d}=\sqrt{1+16-1}=4$.
		\end{itemchoice}
	}
\end{ex}

%G:\My Drive\CODE12-2024\DE-ON-THEO BAI\2H5-TACH DE\Bai4-De2.tex
\begin{ex}%[2H5H3-2]
	Trong không gian $Oxyz$, cho ba điểm $A(1; 2; -4)$, $B(1; -3; 1)$, $C(2; 2; 3)$.
	\choiceTF
	{\True Bán kính của mặt cầu $(S_4)$ đi qua ba điểm $A$, $B$, $C$  và có tâm nằm trên mặt phẳng $(Oxy)$ là $R = \sqrt{26}$}
	{Mặt cầu $\left(S_1 \right)$ tâm $A$, bán kính $R=1$ có phương trình là: $\left(x-1\right)^2+\left(y-2\right)^2+\left(z-4\right)^2=1$}
	{\True Bán kính của mặt cầu $\left(S_2 \right)$ có tâm là $A$ và đi qua điểm $C$ là $\sqrt{50}$}
	{\True Mặt cầu $\left(S_3 \right)$ nhận $AB$ làm đường kính có phương trình là: $\left(x-1\right)^2+\left(y+\dfrac{1}{2} \right)^2+\left(z+\dfrac{3}{2} \right)^2=\dfrac{25}{2}$}
	\loigiai{
		\begin{itemchoice}
			\itemch \textbf{Đúng}.\\
			Gọi phương trình mặt cầu $(S_4)$ có dạng $x^2+y^2+z^2-2ax-2by-2cz+d=0$, với tâm $I(a; b; c)$.\\
			Ta có $I(a;b;c)\in (Oxy) \Rightarrow c=0$. \\
			Vì $ \heva{&{A\in (S)} \\&{B\in (S)} \\&{C\in (S)}}
			\Rightarrow	 \heva{&{-2a-4b+d=-21} \\ &{-2a+6b+d=-11} \\&{-4a-4b+d=-17}}
			\Leftrightarrow \heva{&{a=-2} \\&{b=1} \\&{d=-21}}$\\
			Suy ra $R=\sqrt{a^2+b^2+c^2-d}=\sqrt{4+1+0+21}=\sqrt{26}$.
			\itemch \textbf{Sai}.\\
			Mặt cầu $\left(S_1 \right)$ tâm $A$, bán kính $R=1$ có phương trình là $\left(x-1\right)^2+\left(y-2\right)^2+\left(z+4\right)^2=1$.
			\itemch \textbf{Đúng}.\\
			Ta có $\overrightarrow{AC}=\left(1;0;7\right) \Rightarrow AC=\sqrt{50}$.\\
			Suy ra bán kính của mặt cầu $\left(S_2 \right)$ có tâm là $A$ và đi qua điểm $C$ là $AC=\sqrt{50}$.
			\itemch \textbf{Đúng}.\\
			Ta có $\overrightarrow{AB}=\left(0;-5;5\right) \Rightarrow AB=5\sqrt{2}$.\\
			Mặt cầu $\left(S_3 \right)$ nhận $AB$ làm đường kính nên có tâm $I\left(1;-\dfrac{1}{2};-\dfrac{3}{2} \right)$, $R=\dfrac{AB}{2}=\dfrac{5\sqrt{2}}{2}$, có phương trình là
			\[
			\left(x-1\right)^2+\left(y+\dfrac{1}{2} \right)^2+\left(z+\dfrac{3}{2} \right)^2=\dfrac{25}{2}.
			\]
		\end{itemchoice}
	}
\end{ex}
\Closesolutionfile{ans}

\subsection*{\indam{PHẦN III. Câu trắc nghiệm trả lời ngắn. Thí sinh trả lời từ câu 1 đến câu 6 vào ô kết quả.}}
	\setcounter{ex}{0}
	\Opensolutionfile{ans}[ans/B4-De2-3]

\begin{ex}%[2H5H3-4]
	Trong không gian $Oxyz$, cho phương trình $x^2+y^2+z^2-4x+2my+3m^2-2m=0$ với $m$ là tham số. Tính tổng tất cả các giá trị nguyên của $m$ để phương trình đã cho là phương trình mặt cầu.\\
	\shortans[oly]{$1$}
	\loigiai{
		Phương trình  $x^2+y^2+z^2-4x+2my+3m^2-2m=0$ là phương trình mặt cầu khi và chỉ khi 
		\[
		-2m^2+2m+4> 0\Leftrightarrow m\in \left(-1; 2\right).
		\]
		Do $m\in \mathbb{Z}\Rightarrow m\in \left\{0; 1\right\}$.
		Vậy tổng tất cả các giá trị nguyên của $m$ bằng $1$.
	}
\end{ex}

\begin{ex}%[2H5V3-4]
	Trong không gian $Oxyz$ (đơn vị của các trục tọa độ là ki--lô-mét), đài kiểm soát không lưu sân bay có tọa độ $\left(-64;128;64\right)$. Máy bay bay trong phạm vi cách đài kiểm soát $500$ km thì sẽ hiển thị trên màn hình ra đa. Một máy bay $N$ xuất hiện trên màn hình ra đa và một máy bay $M$ nằm trong mặt phẳng $(P)\colon x-2y+2z-1458=0$ sao cho hai máy bay $M$, $N$ thuộc đường thẳng có vectơ chỉ phương là $\overrightarrow{u}=\left(1;1;1\right)$. Khoảng cách nhỏ nhất giữa hai máy bay $M$, $N$ là bao nhiêu km? (kết quả làm tròn đến hàng đơn vị)\\
	\shortans[oly]{$260$}
	\loigiai{
		\immini{
			Máy bay $N$ xuất hiện trên màn hình ra đa nên $N$ thuộc mặt cầu $(S)$ có tâm $I\left(-64;128;64\right)$, bán kính $R=500$.\\
			Ta có 
			\[	
			\mathrm{d}\left(I,(P)\right)=\dfrac{\left|x_I-2y_I+2z_I-1458\right|}{\sqrt{1^2+2^2+2^2}}=550> R
			\] nên $(P)$ và $(S)$ không giao nhau.\\	
			Gọi $\alpha$ là góc tạo bởi $MN$ và mặt phẳng $(P)$.\\
			Gọi $H$ là hình chiếu của $N$ lên mặt phẳng $(P)$.
		}{\begin{tikzpicture}[scale=1, font=\footnotesize,line join=round, line cap=round, >=stealth, every node/.style={scale=0.8}]
				\coordinate (I) at (0,0);
				\coordinate (N) at (0,-1);
				\coordinate (H) at (0,-3);
				\coordinate (M) at (-1,-3);
				\path let \p1=(I),\p2=(N),\n1={veclen(\x2-\x1,\y2-\y1)} in \pgfextra{\xdef\Temp{\n1}};
				\draw (I) circle (\Temp);
				\coordinate (A) at (-2,-3);
				\coordinate (B) at (4,-3);
				\coordinate (C) at (1.5,-2);
				\coordinate (D) at ($(N)+(C)-(M)$);
				\draw[->] (C)--(D) node[midway,right]{$\vec{u}$};
				\draw(A)--(B) (I)--(H) (M)--(N);
				\foreach \i/\g in {M/-90,H/-90,N/-50,I/180} \fill (\i) circle (1.2pt) node[shift={(\g:3mm)},scale=0.8]{$\i$};
			\end{tikzpicture}
		}
		\noindent
		Mặt khác $\overrightarrow{MN}$ cùng phương với véc-tơ $\overrightarrow{u}=\left(1;1;1\right)$ và $\overrightarrow{n_P}=(1;-2;2)$ suy ra 
		\[
		\sin \alpha=\dfrac{\left|\overrightarrow{u}\cdot \overrightarrow{n_P}\right|}{\left|\overrightarrow{u}\right|\cdot \left|\overrightarrow{n_P}\right|}=\dfrac{1}{3\sqrt{3}}.
		\]
		Khi đó $MN=\dfrac{NH}{\sin \alpha} \ge \dfrac{NH_{\min}}{\sin \alpha}=\dfrac{\mathrm{d}\left(I,(P)\right)-R}{\sin \alpha}=\dfrac{550-500}{\dfrac{1}{3\sqrt{3}}}=150\sqrt{3} \approx 260$.
	}
\end{ex}

\begin{ex}%[2H5V3-2]
	Một vỏ kem ốc quế là một loại bánh khô, hình nón $(N)$ trong không gian $Oxyz$, thường được làm bằng một chiếc bánh xốp dùng để đặt kem vào và cầm ăn mà không cần bát hoặc muỗng. Người ta thả vào vỏ kem $(N)$ một viên kem vani hình cầu có đính hai viên socola nhỏ tại hai vị trí $A(2; 1; 3)$ và $B(6; 5; 5)$ sao cho đường kính $AB$ có $B$ là tâm đường tròn đáy khối nón. Khi thể tích của khối nón $(N)$ nhỏ nhất thì mặt phẳng qua đỉnh $S$ của khối nón $(N)$ và song song với mặt phẳng chứa đường tròn đáy của $(N)$ có phương trình $2x+by+cz+d=0$. Tính giá trị của biểu thức $T=b+c+d$.\\
	\shortans[oly]{$12$}
	\loigiai{
		\immini{
			Gọi chiều cao khối nón $SB=h \left(h > 0\right)$ và bán kính đường tròn đáy $BC=R$.\\
			Ta có $V=\dfrac{1}{3} \pi R^2 h $. \hfill (1)\\
			Ta có $	\overrightarrow{AB}=(4; 4; 2)\Rightarrow AB=6$.\\
			Xét mặt cầu có đường kính $AB$.\\
			Bán kính là $r=\dfrac{AB}{2}=3$ và tâm $I(4; 3; 4)$.\\
			Vì $\triangle SHI$ đồng dạng với $\triangle SBC$ nên
			\[ 
			\dfrac{SI}{SC}=\dfrac{IH}{BC} \Leftrightarrow \dfrac{h-3}{\sqrt{h^2+R^2}}=\dfrac{3}{R}.
			\]
			\[\Leftrightarrow \dfrac{\left(h-3\right)^2}{h^2+R^2}=\dfrac{9}{R^2} \Leftrightarrow R^2\left[\left(h-3\right)^2-9\right]=9h^2\Leftrightarrow R^2=\dfrac{9h^2}{h^2-6h} \tag{2}.
			\]
			Thay $(2)$ vào $(1)$ ta có\\
			$V=\dfrac{1}{3} \pi\cdot \dfrac{9h^2}{h^2-6h}\cdot h=3\pi\cdot \dfrac{h^2}{h-6}$ với $h > 6$.
		}{
			\begin{tikzpicture}[scale=0.8, font=\footnotesize,line join=round, line cap=round, >=stealth, every node/.style={scale=0.8}]
				\def\a{3}
				\def\b{1.0}
				\pgfmathsetmacro\h{\a*sqrt(3)}
				\pgfmathsetmacro\r{\h/3}
				\pgfmathsetmacro\g{asin(\b/\h)}
				\pgfmathsetmacro\xo{\a *cos(\g)}
				\pgfmathsetmacro\yo{\b *sin(\g)}
				\begin{scope}[rotate = 180]
					\path 
					(\xo,\yo) coordinate (Mr) (0,0) coordinate (B) (90:\h) coordinate (S) (0:\a) coordinate (C)  (-\xo,\yo) coordinate (Ml) ($(B)!1/3!(S)$) coordinate (I) ($(B)!2!(I)$) coordinate (A);
					\coordinate (H) at ($(S)!(I)!(C)$);
					\draw[dash pattern=on 2pt off 1.5pt] (S)--(B) (I)--(H)
					(I) circle (\r) (I) ellipse ({\r} and {\r/3});
					\draw (Mr) arc (\g:180-\g:{\a} and {\b}) (S)--(Ml) arc(180-\g:360+\g:{\a} and {\b})--cycle (B)--(C);
					\foreach \x /\gN in {B/-90,S/90,C/0,I/180,A/-40,H/60}
					\fill[black] (\x) circle (1.2pt)($(\x)+(\gN:4mm)$) node {$\x$};
				\end{scope}
			\end{tikzpicture}
		}
		Xét $V'=3\pi\cdot \dfrac{2h\left(h-6\right)-h^2}{\left(h-6\right)^2}=3\pi\cdot \dfrac{h^2-12h}{\left(h-6\right)^2}$.\\
		Ta được BBT như sau:
		\begin{center}
			\begin{tikzpicture}[scale=0.8, font=\footnotesize,line join=round, line cap=round, >=stealth, every node/.style={scale=0.8}]
				\tkzTabInit[nocadre=true,lgt=1.2,espcl=2.5,deltacl=0.6]
				{$h$ /0.6,$V'$ /0.6,$V$ /2.5}
				{$-\infty$,$0$,$6$,$12$,$+\infty$}
				\tkzTabLine{,+,$0$,-,d,-,$0$,+,}
				\tkzTabVar{-/$-\infty$,+/$V(0)$,-D+/$-\infty$/$+\infty$,-/$V(12)$,+/$+\infty$}
				
			\end{tikzpicture}
		\end{center}
		Do $h > 6$ nên $V_{\min}$ khi $SB=h=12$ $\Rightarrow A$ là trung điểm của $SB$ $\Rightarrow S(-2; -3; 1)$.\\
		Khi đó $(P)$ đi qua $S$, vuông góc với $AB$ nên có một VTPT $\overrightarrow{n}=\overrightarrow{AB}=(4; 4; 2)$ hay $\overrightarrow{n}=(2; 2;1)$.\\
		Suy ra $(P)\colon  2(x+2)+2(y+3)+z-1=0\Leftrightarrow (P)\colon 2x+2y+z+9=0$.\\
		Vậy $T=b+c+d=2+1+9=12$.
	}
\end{ex}

\begin{ex}%[2H5V3-2]
	Trong không gian $Oxyz$, cho mặt cầu $(S)\colon x^2+y^2+z^2-2x-4y+6z-13=0$ và đường thẳng $d\colon \heva{& x=-1+t \\ & y=-2+t \\ & z=1+t}$. Gọi $M(a; b; c)$ với $a < 0$ là điểm thuộc đường thẳng $d$ sao cho từ $M$ kẻ được ba tiếp tuyến $MA$, $MB$, $MC$ đến mặc cầu $(S)$ ($A,B,C$ là các tiếp điểm) thỏa mãn $\widehat{AMB}=60^{\circ}$; $\widehat{BMC}=90^{\circ}$; $\widehat{CMA}=120^{\circ}$. Tính giá trị của biểu thức $P=a+b+c$.\\
	\shortans[oly]{$-2$}
	\loigiai{
		\immini{
			Ta có mặt cầu $(S)$ có tâm $I\left(1;2;-3\right)$, bán kính $R=3\sqrt{3}$.\\
			Đặt $MA=MB=MC=a$.\\
			Tam giác $MAB$ đều nên $AB=a$.\\
			Tam giác $MBC$ vuông cân tại $M$ nên $BC=a\sqrt{2}$.\\
			Tam giác $MCA$ có $\widehat{CMA}=120^{\circ}$ nên $CA=a\sqrt{3}$.\\
			Xét tam giác $ABC$ có $AB^2+BC^2=AC^2$ nên tam giác $ABC$ vuông tại $B$ hay tam giác $ABC$ nội tiếp đường tròn đường kính $AC$ và $AH=\dfrac{1}{2} AC=\dfrac{a\sqrt{3}}{2}$.
		}{
			\begin{tikzpicture}[scale=0.7, line join=round, line cap=round,>=stealth]
				\coordinate (H) at (0,0);
				\coordinate (I) at (0,-3);
				\coordinate (M) at (0,5);
				%\draw (H) ellipse (4 cm and 1.5 cm);
				\coordinate (A) at ($(H) + (110:4 and 1.5)$);
				\coordinate (C) at ($(H) + (-30:4 and 1.5)$);
				\coordinate (B) at ($(H) + (-150:4 and 1.5)$);
				\coordinate (E) at ($(H) + (-90:4 and 1.5)$);
				\draw pic[draw,angle radius=0.3cm]{right angle=I--B--M};
				\draw pic[draw,angle radius=0.3cm]{right angle=I--A--M};
				\draw pic[draw,angle radius=0.3cm]{right angle=M--C--I};
				\draw  (B)--(M)--(C) (I)--(B) (I)--(C) (B)--(C);
				\draw[dashed] (A)--(B) (C)--(A) (H)--(A) (H)--(B) (H)--(C) (A)--(M)--(H) (I)--(A) (H)--(I);
				\foreach \i/\g in {A/130,B/-130,C/0,M/90,I/-90,H/45}
				\fill[black] (\i) circle(1pt)+(\g:3mm)node[scale=0.7]{$\i$};
			\end{tikzpicture}
		}	
		\noindent
		Xét tam giác vuông $IAM$ có
		\[
		\dfrac{1}{AH^2}=\dfrac{1}{AM^2}+\dfrac{1}{AI^2} \Rightarrow MA=a=3\Rightarrow IM^2=AM^2+AI^2=36.
		\]
		Mặt khác $M\in (d)\Rightarrow M\left(-1+t;-2+t;1+t\right)$.\\
		Hay $\left(t-2\right)^2+\left(t-4\right)^2+\left(t+4\right)^2=36$ $\Leftrightarrow 3t^2-4t=0\Leftrightarrow \hoac{ & t=0 \\ & t=\dfrac{4}{3}.}$\\
		Suy ra $M\left(-1;-2;1\right)$ hay $M\left(\dfrac{1}{3};-\dfrac{2}{3};\dfrac{7}{3} \right)$.\\
		Vì giả thiết cho $a < 0$ nên $a=-1;b=-2;c=1$. Khi đó $P=a+b+c=-2$.
	}
\end{ex}

\begin{ex}%[2H5V3-4]
	Trong không gian $Oxyz$ (đơn vị của các trục tọa độ là ki--lô-mét), một trạm thu phát sóng điện thoại di động có đầu thu đặt tại điểm $I\left(1;2;2\right)$ biết rằng bán kính phủ sóng của trạm là $3$ km. Hai người sử dụng điện thoại lần lượt tại $M\left(4;-4;2\right)$ và $N\left(6;0;6\right)$. Gọi $E(a; b; c)$ với $a < 0$ là một điểm thuộc ranh giới vùng phủ sóng của trạm sao cho tổng khoảng cách từ $E$ đến vị trí $M$ và $N$ lớn nhất. Tính $T=a+b+c$.\\
	\shortans[oly]{$4$}
	\loigiai{
		\immini{
			Xét mặt cầu $(S)$ có tâm $I\left(1;2;2\right)$, bán kính $R=3$.\\
			Ta có $IM=IN=3\sqrt{3} > R$ nên cả hai điểm $M$, $N$ đều nằm ngoài mặt cầu $(S)$.\\
			Gọi $H$ là trung điểm của $MN$ suy ra $H\left(5;-2;4\right)$ và 
			\[
			EH^2=\dfrac{EM^2+EN^2}{2}-\dfrac{MN^2}{4}.
			\]
			Ta có $\left(EM+EN\right)^2\le 2\left(EM^2+EN^2\right)=2\left(EH^2+\dfrac{MN^2}{4} \right)$.
		}{
			\begin{tikzpicture}[scale=0.6, font=\footnotesize,line join=round, line cap=round, >=stealth, every node/.style={scale=0.8}]
				\def\r{2.0}
				\def\x{2.0}
				\def\y{0.8}
				\path
				(0,0) coordinate (I)
				(-4,0) coordinate (M)
				(-\r,0) coordinate (A)
				(\r,0) coordinate (B)
				($(I)+({\x*cos(-120)},{\y*sin(-120)})$) coordinate (E)
				($(I)+({2.5*\x*cos(-120)},{2.5*\y*sin(-120)})$) coordinate (H)
				($(E)!2!(I)$) coordinate (E')
				($(M)!2!(H)$) coordinate (N)
				;
				\draw (I) circle(\r)
				(B) arc (0:-180:\x cm and \y cm)
				(E)--(H) (M)--(N)
				;
				\draw[dashed] (B) arc (0:180:\x cm and \y cm)
				(I)--(B)node[below,midway]{$R$} (E')--(I)--(E);
				\foreach \i/\g in {I/90,E/-90,E'/90,H/-120,M/90,N/-90}
				\fill[black] (\i) circle(1pt)+(\g:3mm)node[scale=1]{$\i$};
			\end{tikzpicture}
		}
		\noindent
		Khi đó $\left(EM+EN\right)$ lớn nhất khi $EH$ lớn nhất khi và chỉ khi $E$ là giao điểm của $IH$ và mặt cầu $(S)$.\\
		$IH$ có phương trình là $\heva{& x=1+2t \\ & y=2-2t \\ & z=2+t.}$ \hfill (1)\\
		Mặt cầu $(S)$ có phương trình $\left(x-1\right)^2+\left(y-2\right)^2+\left(z-2\right)^2=9.$ \hfill (2)\\
		Thay (1) vào (2), ta được
		\[(2t)^2+\left(-2t\right)^2+t^2=9\Leftrightarrow t=\pm 1.
		\]
		Suy ra có hai điểm $E'\left(3;0;3\right)$ hoặc $E\left(-1;4;1\right)$. \\
		Vì $a < 0$ nên $E\left(-1;4;1\right)$ hay $T=a+b+c=4$.
	}
\end{ex}

\begin{ex}%[2H5H3-4]
	\immini[thm]{Người ta muốn thiết kế một bồn chứa khí hoá lỏng hình cầu bằng phần mềm 3D (tham khảo hình vẽ). Cho biết phương trình bề mặt của bồn chứa là $(S)\colon  \left(x-2\right)^2+y^2+\left(z+1\right)^2=1$. Phương trình mặt phẳng chứa nắp là $(P)\colon  z-6=0$. Tính khoảng cách từ tâm bồn chứa đến mặt phẳng chứa nắp.\\
	\shortans[oly]{$7$}
}{
\includegraphics[scale=0.7]{image/bonkhi}}
	\loigiai{
		Tâm của bồn chứa $I\left(2;0;-1\right)$.\\
		Khoảng cách từ tâm bồn chứa đến mặt phẳng chứa nắp là
		\[
		\mathrm{d}\left(I,(P)\right)=\dfrac{\left|-1-6\right|}{\sqrt{1^2}}=7.
		\]
	}
\end{ex}

\centerline{---HẾT---}
\Closesolutionfile{ans}
%\newpage
%%=====================
%\begin{center}
%\textbf{\large BẢNG ĐÁP ÁN}
%\end{center}
%\noindent\textbf{ĐÁP ÁN PHẦN I}
%\inputansbox{10}{ans/B4-De2-1}
	
%\noindent\textbf{ĐÁP ÁN PHẦN II}
%\inputansbox[2]{2}{ans/B4-De2-2}
	
%\noindent\textbf{ĐÁP ÁN PHẦN III}
%\inputansbox[3]{6}{ans/B4-De2-3}


